% Chapter 4 — local volatility.
% Sections live in this directory; figures and macros regenerate via
% scripts/ch04/gen_figures.py from data/ (frozen snapshot + deterministic
% synthetics).  Notation additions are recorded in NOTATION.md.
\renewcommand{\figdir}{figures/ch04}
\chapter[Local volatility]{Local Volatility}
\label{ch:lv}

\chaptersynopsis{Chapters~2 and~3 fitted one law per expiry and let a
calendar condition relate neighbors.  This chapter asks for more: one
diffusion whose marginal laws are all of the surface's laws at once.  The
Dupire equation supplies it, and in the volatility chart its coefficient is
a ratio of the two objects the static chapters policed --- the calendar
increment over the Durrleman factor --- so the field exists exactly where
the surface is valid.  Read as a formula, the ratio divides two derivatives
of noisy data and fails structurally; read as a pricing map, it poses a
well-behaved inverse problem: a piecewise-affine sheet whose positivity is
a box constraint on finitely many numbers, priced by a monotone implicit
march whose step matrix is an M-matrix, and differentiated exactly behind
the same factorization.  The quotes do not determine the sheet --- two
surfaces ten volatility points apart reprice the same quote set within a
fraction of a spread --- so calibration returns a declared representative
of an equivalence class, and what the discretization leaks is measured,
never assumed.}

\section{One diffusion for the whole surface}
\label{sec:lvfield}

The static chapters treated an expiry as a law.  Chapter~2 built the law
first and read the smile off it; Chapter~3 worked in the volatility chart
and policed validity there.  Either way, a fitted surface is at this point a
\emph{family} of laws, one per expiry, related pairwise by the calendar
comparison of \cref{sec:calendar}: at equal log-moneyness, later normalized
calls dominate earlier ones, which is convex order of the normalized
returns.  Convex order is exactly the condition under which the family is
jointly realizable: by Kellerer's theorem \citep{Kellerer1972}, \emph{some}
martingale has these marginals.  But Kellerer's theorem is an existence
statement.  It does not say which martingale, it does not give the process
any structure a computation could use, and it leaves open what additional
information a particular choice smuggles in.

Dupire's observation \citep{Dupire1994} resolves all three questions inside
one class.  Adopt, as a standing hypothesis, the deterministic-carry
setting of \cref{sec:calendar}, under which $Y_\tau=S_\tau/F_{0,\tau}$ is a
single positive martingale whose time-$\tau$ marginal is Chapter~2's
normalized return for the expiry-$\tau$ slice.  Consider diffusions of the
form
\begin{equation}
\dd Y_\tau=\sqrt{v(\tau,Y_\tau)}\;Y_\tau\,\dd W_\tau,\qquad Y_0=1,
\label{eq:lvsde}
\end{equation}
where $v(\tau,y)\ge0$ is a deterministic field, the \emph{local
variance}; its square root $\sigma_{\rm loc}=\sqrt v$ is the local
volatility.  Within
this class the correspondence between coefficient and marginals is exact in
both directions: the field determines every marginal law by solving a
forward equation, and the marginals determine the field by a pointwise
formula, wherever the density is positive.  A valid surface is therefore
not merely \emph{consistent with} some model --- it \emph{is} the
coefficient field of one canonical model, written in different coordinates.
This chapter is about that change of coordinates: how to compute it, why
the obvious way of computing it fails on real data, and what a market's
finitely many quotes do and do not say about the field.

Two remarks bound the object's ambitions before we start.  First, the
diffusion \eqref{eq:lvsde} is a \emph{representative}, not a discovery.  By
Gy\"ongy's theorem \citep{Gyongy1986}, any It\^o martingale (under mild
regularity) with the same marginals is mimicked by \eqref{eq:lvsde} with
$v(\tau,y)$ equal to the conditional expectation of its instantaneous
variance given $Y_\tau=y$: the local-volatility model is the Markovian
projection common to every dynamics the statics allow.  Statements that
depend only on marginal laws --- every European price, every quantity of
Chapters~2--3 --- are safe; statements about joint laws, such as how the
smile moves when the spot moves, are not determined by the surface at all
and are postponed to Chapter~9.  Second, the class \eqref{eq:lvsde}
contains no jumps.  A genuine discontinuity --- an earnings gap, a
scheduled announcement --- can only be imitated by a burst of diffusive
variance; scheduled events belong to a change of clock, which is
Chapter~8's subject.  Throughout this chapter $\tau$ is the
variance time used since Chapter~2, $y=e^{k}$ the normalized strike, and
$c(\tau,y)$ the normalized call: rates, borrow and dividends live inside
the forward (Chapter~6), and the field $v$ is variance per unit variance
time.

The plan follows the two directions of the correspondence.
\Cref{sec:lveq} derives the forward equation, shows that positivity of the
field propagates both static no-arbitrage families at once, and identifies
the inverse formula as a ratio of two objects the reader already owns.
\Cref{sec:lvillposed} demonstrates why that formula must not be applied to
data.  \Cref{sec:lvsheet,sec:lvlattice} build the objects that replace it:
a finite-dimensional field whose validity is a box constraint, and a
discretization that inherits the continuous theorem instead of merely
approximating it.  \Cref{sec:lvinfo} asks what the quotes actually
determine, and \cref{sec:lvcalib,sec:lvexamples} assemble the calibration
and run it on the book's frozen data.

\section{The equation the surface must obey}
\label{sec:lveq}

\subsection{Forward: from the field to every marginal}

Fix a bounded measurable field $v\ge0$ and let $Y$ solve \eqref{eq:lvsde}.
Everything follows from watching the call price evolve in maturity.

\begin{proposition}[The forward Dupire equation]
\label{prop:lvforward}
Suppose $Y_\tau$ has a density $f(\tau,\cdot)$ such that
$v y^{2}f$ and $y\,\partial_y(v y^{2}f)$ are integrable and vanish at
$y=0$ and as $y\to\infty$.
Then the normalized call $c(\tau,y)=\E[(Y_\tau-y)^+]$ satisfies
\begin{equation}
\boxed{\;\partial_\tau c(\tau,y)
=\tfrac12\,v(\tau,y)\,y^{2}\,\partial_{yy}c(\tau,y),
\qquad c(0,y)=(1-y)^{+}.\;}
\label{eq:lvdupire}
\end{equation}
\end{proposition}

\begin{proof}
By It\^o's formula, the density satisfies the Fokker--Planck equation
$\partial_\tau f=\tfrac12\,\partial_{yy}\!\big(v y^{2}f\big)$ in the weak
sense.  Integrate it against the payoff:
\[
\partial_\tau c(\tau,y_0)
=\int_0^\infty (y-y_0)^{+}\,\partial_\tau f(\tau,y)\,\dd y
=\tfrac12\int_0^\infty (y-y_0)^{+}\,
  \partial_{yy}\big(v y^{2}f\big)\,\dd y .
\]
Integrating by parts twice moves the two derivatives onto the payoff,
whose derivatives in $y$ are the step $\1_{\{y>y_0\}}$ and then the point
mass $\delta_{y_0}$, while the boundary terms vanish by the decay
hypothesis.  What remains is
$\tfrac12\,v(\tau,y_0)\,y_0^{2}f(\tau,y_0)$, and
$f(\tau,y_0)=\partial_{yy}c(\tau,y_0)$ is the density identity of
\cref{sec:intro}.  The initial condition is $Y_0=1$.
\end{proof}

Equation \eqref{eq:lvdupire} is a \emph{forward} parabolic equation in the
strike variable with the payoff as initial data: heat flow, with
$\tfrac12 v y^{2}$ as a spatially varying conductivity.  One field, one
linear march, and every marginal law of the surface comes out at once ---
this is the computational engine of the whole chapter.  The analogy with
heat is not decoration; it is the source of every structural property we
will use.  Heat flow spreads bumps and never sharpens them, and spreading,
translated into option language, is exactly the no-arbitrage direction:

\begin{proposition}[Positive local variance propagates no-arbitrage]
\label{prop:lvpropagate}
Let $c$ solve \eqref{eq:lvdupire} with $v\ge0$ bounded.  Then for every
$\tau>0$: (i) $\partial_{yy}c(\tau,\cdot)\ge0$ --- each slice is
butterfly-free; (ii) $\partial_\tau c(\tau,y)\ge0$ --- the normalized call
is non-decreasing in maturity at fixed $y$, which is the calendar order of
\cref{sec:calendar}.
\end{proposition}

\begin{proof}[Proof sketch]
Differentiating \eqref{eq:lvdupire} twice in $y$ shows that
$f=\partial_{yy}c$ satisfies
$\partial_\tau f=\tfrac12\,\partial_{yy}(v y^{2}f)$ --- the
Fokker--Planck equation again --- with initial datum
$\partial_{yy}(1-y)^{+}=\delta_{1}\ge0$.  Such equations preserve
nonnegativity: probabilistically, $f(\tau,\cdot)$ \emph{is} the density of
the diffusion \eqref{eq:lvsde} started at $1$; analytically, apply the
parabolic maximum principle to the adjoint equation.  With $f\ge0$
established, (ii) is \eqref{eq:lvdupire} itself:
$\partial_\tau c=\tfrac12 v y^{2}f\ge0$.
\end{proof}

Pause on what this proposition buys, because it is the entire strategy of
the chapter.  In Chapters~2 and~3, validity was work: the LQD chart made
butterfly freedom an identity within one slice, the volatility chart had to
police $g_{\rm D}\ge0$, and calendar order across slices was in both cases
a separate condition.  Here, positivity of \emph{one} function of two
variables propagates \emph{both} families across the whole surface, for
free, forever.  The rest of the chapter is the discipline of not
squandering that theorem: choose a parametrization in which $v>0$ is
trivial to enforce (\cref{sec:lvsheet}), choose a discretization that
inherits the maximum principle rather than merely approximating it
(\cref{sec:lvlattice}), and then measure what little the finite grid leaks
instead of assuming the theorem carried over.

\subsection{Backward: from the surface to the field}

Read \eqref{eq:lvdupire} the other way.  If the surface is given, the field
is determined pointwise,
\begin{equation}
v(\tau,y)=\frac{\partial_\tau c(\tau,y)}
{\tfrac12\,y^{2}\,\partial_{yy}c(\tau,y)},
\label{eq:lvextract}
\end{equation}
wherever the density in the denominator is positive.  This is the promised
uniqueness: within the class \eqref{eq:lvsde}, at most one field is
consistent with a given surface, and \eqref{eq:lvextract} computes it.  The
same formula becomes most instructive in the volatility chart.

\begin{proposition}[The field in the volatility chart]
\label{prop:lvratio}
Write the surface in implied form, $c=\Black(k,w(k,\tau))$ with
$w(\cdot,\tau)$ the total-variance slice.  Then \eqref{eq:lvextract} reads
\begin{equation}
v(\tau,e^{k})=\frac{\partial_\tau w(k,\tau)}{g_{\rm D}(k)},
\label{eq:lvratio}
\end{equation}
with $g_{\rm D}$ the Durrleman factor of \eqref{eq:durrleman}, evaluated on
the slice at maturity $\tau$ \citep{Gatheral2006}.
\end{proposition}

\begin{proof}
Numerator: at fixed $k$, $\partial_\tau c=\partial_w\Black\cdot
\partial_\tau w$.  With $d_\pm=-k/\sqrt w\pm\sqrt w/2$ one has
$d_+^{2}-d_-^{2}=(d_++d_-)(d_+-d_-)=-2k$, hence
$\phin(d_+)=e^{k}\phin(d_-)$, and
\[
\partial_w\Black
=\phin(d_+)\,\partial_w d_+-e^{k}\phin(d_-)\,\partial_w d_-
=\phin(d_+)\,\partial_w(d_+-d_-)
=\frac{e^{k}\phin(d_-)}{2\sqrt w}.
\]
Denominator: $\partial_{yy}c=f(\tau,y)=f_X(k)/y$ by the change of variable
$y=e^{k}$, and \eqref{eq:durrleman} gives
$f_X(k)=g_{\rm D}(k)\,\phin(d_-)/\sqrt w$.  Substituting both into
\eqref{eq:lvextract},
\[
v=\frac{e^{k}\phin(d_-)\,\partial_\tau w/\sqrt w}
       {y^{2}\,f_X(k)/y}
 =\frac{\phin(d_-)\,\partial_\tau w}{\sqrt w\,f_X(k)}
 =\frac{\partial_\tau w}{g_{\rm D}(k)}.\qedhere
\]
\end{proof}

The reader has met both halves of this ratio before, as the two families of
static validity.  The numerator is the calendar increment: nonnegative
exactly when adjacent slices are ordered in the sense of
\cref{sec:calendar}.  The denominator is positive density written in the
volatility chart, the object Chapter~3 policed slice by slice.  So
\eqref{eq:lvratio} is a compatibility theorem disguised as a formula:
\emph{a local-volatility field exists precisely where the surface is
valid}, and each mode of static arbitrage corresponds to its own sign
failure --- a calendar violation makes the numerator negative, a butterfly
violation the denominator.  This chapter therefore takes as input
hypothesis what Chapters~2--3 produced as output: slices that are
butterfly-free and calendar-ordered, in the normalized units of
\cref{sec:calendar}.  Where the hypothesis fails, no field exists and
\eqref{eq:lvratio} says so by changing sign.

\begin{figure}[htbp]
\centering
\paperfig{\textwidth}{fig_lv_ratio.pdf}
\caption{The field as a ratio, on the frozen SPY surface (stored per-expiry
fits of the snapshot, total variance interpolated linearly in $\tau$
between the eight expiries; horizontal bars mark each expiry's quoted
span).  (a)~The calendar increment $\partial_\tau w$: piecewise constant in
$\tau$ by construction, steepening toward the short-dated put side.
(b)~The Durrleman factor of each slice.  (c)~Their ratio, as local
volatility: the put wall, the quiet call side, and the short-dated
structure.  Grey cells are inadmissible (a negative numerator or
denominator): they amount to \MacLvRatioBadSharePct\% of the displayed
mesh, all in extrapolated territory --- restricted to the common quoted
support of each expiry pair the share is
\MacLvRatioBadQuotedSharePct\%, and the extracted local volatility there
ranges from \MacLvRatioLocMinPct\% to \MacLvRatioLocMaxPct\%.  Validity
where the data speaks, sign failures where conventions extrapolate: the
ratio is already a diagnostic instrument.}
\label{fig:lvratio}
\end{figure}

\Cref{fig:lvratio} evaluates the ratio on the book's frozen SPY surface ---
the eight fitted slices of the snapshot, interpolated linearly in total
variance between expiries.  The exercise makes three points at once.  The
field inherits recognizable structure from the smile: a steep put-side wall,
a quiet call side, short-dated sharpness.  The two sign conditions are
visibly the binding ones: every inadmissible cell in the picture sits
outside the common quoted support of the expiries that bracket it, in
territory owned by extrapolation conventions rather than prices.  And the
ratio was computed here from \emph{fitted, smooth} slices --- an operation
this chapter endorses.  What it must never be applied to is the market's
raw quotes, and the next section shows why.

\section{The wrong direction}
\label{sec:lvillposed}

Formula \eqref{eq:lvextract} invites a procedure: interpolate the quoted
smile, differentiate it twice in strike and once in maturity, divide.  On a
continuum of exact prices this is a theorem.  On a real chain --- a few
dozen noisy quotes per expiry --- it fails, and the failure is structural,
not bad luck.

Quantify it.  Suppose each quoted total variance carries an error of size
$\varepsilon$, and strikes are spaced $\Delta k$ apart.  A second difference
estimates curvature as
\[
\partial_{kk}w\;\approx\;
\frac{w(k-\Delta k)-2\,w(k)+w(k+\Delta k)}{\Delta k^{2}},
\]
and the worst configuration of errors --- alternating signs, which is
exactly what a bid--ask bounce produces --- enters this estimate at size
$4\varepsilon/\Delta k^{2}$.  The noise is amplified by the \emph{inverse square}
of the strike spacing: a finer quote grid makes the estimate worse, without
bound.  There is no stable limit to converge to; numerical differentiation
of noisy data is the canonical ill-posed operation.  And the place this
amplified noise lands could not be worse chosen: by \eqref{eq:lvratio} the
curvature sits inside $g_{\rm D}$, the sign-carrying factor of the
risk-neutral density.  The first
time the noise pushes the estimated density through zero, the extraction
returns a \emph{negative variance} --- not a large error but a
non-object, which no pricing equation, no simulation and no risk system
downstream can absorb.

\begin{figure}[htbp]
\centering
\paperfig{\textwidth}{fig_lv_wrongway.pdf}
\caption{The same synthetic market, read both ways.  Quotes are generated
from a known smooth field (dashed) and carry a deterministic alternating
$\pm$\MacLvWrongRippleBp\,bp volatility ripple --- a stylized bid--ask
bounce.  (a)~The extraction \eqref{eq:lvratio} from the noisy quotes at
$\tau=0.25$.  From 13 quotes per expiry it saw\-tooths between
\MacLvWrongCoarseMinPct\% and \MacLvWrongMaxSpikePct\% around a true ATM
level of \MacLvWrongTruthAtmPct\%; from 31 quotes of the \emph{same}
market --- finer $\Delta k$, same $\pm$\MacLvWrongRippleBp\,bp ---
it fails outright at \MacLvWrongNBad\ strikes (crosses: negative variance
or negative density), exactly the $4\varepsilon/\Delta k^{2}$ scaling at work.
(b)~The same 31-quote noisy chains handed to the forward calibration of
this chapter: an admissible surface by construction, within
\MacLvWrongFwdMaxPts\ volatility points of the truth along the plotted
cross-sections, with reprice residuals of rms
\MacLvWrongQuoteRmsBp\,bp --- the scale of the ripple itself.  The noise
stays in the residual, where it belongs.}
\label{fig:lvwrongway}
\end{figure}

\Cref{fig:lvwrongway} performs the experiment on a synthetic market whose
answer is known (the field and quote recipe are in \cref{app:lvnumerics}).
The left panel is the naive pipeline at two quote densities and needs
little commentary beyond its caption: refining the strike grid, which
improves every well-posed computation, makes this one \emph{worse},
quadratically, until half the strikes report that no diffusion exists.
The right panel previews the alternative developed in the rest of the
chapter.  The same noisy quotes are handed to a calibration that never
differentiates them; every iterate it produces is a genuine positive
field; and the fitted surface tracks the truth through the quoted region
while its residuals absorb the ripple.  One caveat, quantified now and
explained in \cref{sec:lvinfo}: at the edges of the shortest expiry's
quoted span, where a single rippled quote is the only witness, the fitted
field wanders by as much as \MacLvWrongEdgePts\ volatility points.  Noise
rejection is a property of dense data and declared smoothing.

The repair is a change of role, not of formula, and it is due to Dupire's
equation being read in its \emph{forward} direction, as Andreasen and Huge
taught \citep{AndreasenHuge2011}.  Decide that the unknown is the field
$v$ itself.  Parametrize it.  Price it through \eqref{eq:lvdupire} ---
which differentiates only the smooth numerical solution, never the data
--- and ask an optimizer for the positive field whose prices meet the
quotes.  Positivity is imposed on the \emph{parameters}, so every iterate
of the search is an admissible diffusion before any quote is consulted.
Readers of Chapters~2--3 will recognize the design: it is the same
constraint-elimination idea that built the LQD chart and the structural
SVI chart (\cref{sec:charts,sec:svistructural}), one dimension up ---
choose the coordinates in which the constraint set is trivial, and let the
pricing map do the hard work.  Three questions remain, and they organize
the rest of the chapter: what parametrization makes positivity trivial in
two dimensions (\cref{sec:lvsheet}); what discretization of
\eqref{eq:lvdupire} keeps \cref{prop:lvpropagate} alive on a finite grid
(\cref{sec:lvlattice}); and what, exactly, the quotes determine about the
field once the formula \eqref{eq:lvextract} is off the table
(\cref{sec:lvinfo}).

\section{A sheet with finitely many parameters}
\label{sec:lvsheet}

\subsection{Positivity as a box}

The unknown is now a positive function of two variables --- an
infinite-dimensional object with a one-sided constraint, which sounds
harder than anything fitted so far in this book.  The finite version
is built from the simplest ingredient available: straight-line
interpolation.  Picture the smallest example that shows everything.
Take two maturities and three strikes --- a grid of six points --- and
write a positive number on each.  Split each of the two rectangular
cells along a diagonal, giving four triangles.  Inside each triangle,
define the field as the unique plane through its three corner values.
The result is a continuous surface made of flat triangular facets,
passing through the six numbers --- and those six numbers are the
model's only parameters.  The general object just replaces six by a
few hundred.

Two names before the formal definition.  A \emph{tensor grid} is a
rectangular grid: pick a list of maturities and a list of strikes, and
take every combination as a vertex.  The \emph{hat function} of a
vertex is the tent that rises to $1$ there, slopes to $0$ at every
neighboring vertex, and stays $0$ beyond; on each triangle it is the
affine function with those corner values.  Its values at a point are
also called the point's \emph{barycentric weights} --- the three
nonnegative numbers, summing to one, that locate the point inside its
triangle.

\begin{definition}[$P_1$ local-variance sheet]
\label{def:lvsheet}
Fix vertices $(\tau_i,y_j)$, $i=1..n_\tau$, $j=1..n_y$, on a tensor grid,
triangulated by splitting each rectangular cell along one diagonal
(\cref{rem:lvdiag}).  The local variance is the continuous
piecewise-affine interpolant
\begin{equation}
v_\theta(\tau,y)=\sum_{\ell=1}^{n_\theta}\theta_\ell\,
\mathcal H_\ell(\tau,y),\qquad n_\theta=n_\tau n_y,
\label{eq:lvp1}
\end{equation}
where $\mathcal H_\ell$ is the hat function of vertex $\ell$ ---
barycentric on each triangle, equal to $1$ at its own vertex and $0$ at
every other --- and the parameters $\theta_\ell=v(\tau_i,y_j)>0$ are the
\emph{nodal local variances}, indexed flat as $\theta_\ell$ or by row and
column as $\theta_{i,j}$, whichever is convenient.
\end{definition}

With the names in place, read \eqref{eq:lvp1} as plainly as it
deserves: the sum simply interpolates the nodal values.  At any point
of the grid, only the three hats of the surrounding triangle are
nonzero; their values there are the point's barycentric weights; and
$v_\theta$ at that point is the corresponding weighted average of the
three corner parameters.

Why affine pieces, and why triangles?  Because a weighted average with
nonnegative weights cannot leave the range of the numbers it
averages.  That one-line fact is the model's entire validity theory:

\begin{proposition}[Nodal bounds are surface bounds]
\label{prop:lvnodal}
If $v_{\rm lo}\le\theta_\ell\le v_{\rm hi}$ for all $\ell$, then
$v_{\rm lo}\le v_\theta(\tau,y)\le v_{\rm hi}$ everywhere on the
triangulation's hull.
\end{proposition}

\begin{proof}
At any point of the hull, the only nonzero hat functions are the three
belonging to the surrounding triangle, and their values there are
nonnegative and sum to one --- they are the point's barycentric
weights.  Hence
\[
v_\theta(\tau,y)
=\sum_{\ell}\theta_\ell\,\mathcal H_\ell(\tau,y)
\;\le\;v_{\rm hi}\sum_{\ell}\mathcal H_\ell(\tau,y)
=v_{\rm hi},
\]
and the lower bound is the same computation with $v_{\rm lo}$.
\end{proof}

Positivity of a surface has been reduced to positivity of finitely
many numbers: keep every $\theta_\ell$ inside the box
$[v_{\rm lo},v_{\rm hi}]$, and the whole field stays inside it too,
everywhere on the hull --- the region the triangles cover.  A
box is the one constraint that least-squares solvers handle natively
--- clip each coordinate --- with no penalty terms and no feasibility
questions.  This is the same chart construction, one dimension up: as with the
LQD chart of Chapter~2 and the structural chart of
\cref{sec:svistructural}, the coordinates are chosen so that every
point of the parameter set is a valid object, and the optimizer never
learns that a constraint existed.  Combined with
\cref{prop:lvpropagate}, the chain of guarantees is complete before
any data arrives:
\[
\text{box}\;\Longrightarrow\;\text{positive field}
\;\Longrightarrow\;\text{butterfly-free, calendar-ordered surface.}
\]
Every iterate of every fit in this chapter is an admissible diffusion.

\begin{figure}[htbp]
\centering
\paperfig{0.9\textwidth}{fig_lv_sheet.pdf}
\caption{The model, drawn: the calibrated SPY local-volatility sheet of
\cref{sec:lvexamples} (\MacLvSheetRows\ maturity rows $\times$
\MacLvSheetCols\ strike columns, \MacLvSheetVtx\ vertices), with the
triangulation the pricing basis actually uses.  Every dot is one
calibration parameter; every triangle is a region where local variance
is affine in $(\tau,y)$.}
\label{fig:lvsheet}
\end{figure}

\Cref{fig:lvsheet} shows the object this chapter will calibrate: the
SPY sheet of \cref{sec:lvexamples}, with \MacLvSheetRows\ maturity
rows, \MacLvSheetCols\ strike columns, and therefore \MacLvSheetVtx\
vertices, drawn with the very triangulation the pricing uses.  Read it
as a list of numbers, because that is all it is.  Every dot is one
calibration parameter; every triangle is a region where local variance
is an affine function of $(\tau,y)$.  The tall put-side wall, the
sharp structure at the short end, and the flat call-side plain are
nothing but nodal values.  Notice also that the strike columns crowd
toward the money at short maturities; that is deliberate grid design,
and \cref{sec:lvinfo} explains it.

\begin{figure}[htbp]
\centering
\paperfig{\textwidth}{fig_lv_basis.pdf}
\caption{The anatomy of the sheet, on a small demonstration grid.
(a)~The triangulated vertex grid; the shaded star of triangles is the
support of one vertex's hat function.  The cell diagonals are not
uniformly oriented; that is real (\cref{rem:lvdiag}).  (b)~A strike
cross-section along a vertex row: each hat is $1$ at its own vertex,
$0$ at every other, affine between, and they sum to $1$ identically.}
\label{fig:lvbasis}
\end{figure}

\Cref{fig:lvbasis} takes the construction apart on a small
demonstration grid.  In panel~(a), the triangulated vertex grid; the
shaded star of triangles around one vertex is the support of that
vertex's hat function --- the entire region of the surface its
parameter can influence.  A nodal variance is a local dial: turn it,
and only the prices of options whose diffusion passes through that
star can move.  (The cell diagonals in the panel do not all point the
same way; that is real, and \cref{rem:lvdiag} explains it.)  In
panel~(b), a cross-section along one vertex row: each hat rises to $1$
at its own vertex, is $0$ at every other, and the hats sum to $1$
identically --- the partition of unity behind the weighted-average
argument of \cref{prop:lvnodal}.

\begin{remark}[Which diagonal, and does it matter]
\label{rem:lvdiag}
On a tensor grid every cell is a rectangle, and a rectangle's four
corners lie on one circle --- exactly the tie case of the standard
(Delaunay) criterion for well-shaped triangles, which prefers
triangulations in which no vertex falls strictly inside another
triangle's circumscribed circle.  The triangulation routine therefore breaks the tie cell by
cell, deterministically but by no designed rule, and the figures draw
the resulting patchwork faithfully.  The choice is second-order small:
the two candidate interpolants of a cell agree at all four corners and
along all four edges, and differ inside by at most half the cell's
twist,
$\big(\theta_{i,j}+\theta_{i+1,j+1}-\theta_{i,j+1}-\theta_{i+1,j}\big)/2$,
at the center --- zero whenever the four values are locally bilinear.
It is nonetheless a convention, so it is pinned: the triangulation is
computed once, cached, and the same triangles price every iterate.
\end{remark}

\subsection{Beyond the hull: the wing contract}

The grid must end somewhere, and pricing does not.  Below the lowest
strike column, the sheet continues local variance \emph{linearly},
with slope $a$ times the slope of the first interior cell; above the
highest column it is clamped flat.  Three settings of $a$ are
meaningful.  The default, $a=0$, clamps the put wing flat as well.  A
fixed $a>0$ keeps an unquoted put wing rising.  And $a$ is handed to
the optimizer as a free parameter in exactly one circumstance: when a
variance-swap quote is present.  A variance swap pays the realized
variance of the underlying to expiry, and its fair price integrates
the whole smile, deep wings included --- it is the one instrument in
the problem that genuinely reads the deep-put tail, as Chapter~5 will
show.

\begin{remark}[The extrapolation sits outside the positivity guarantee]
\label{rem:lvhull}
\Cref{prop:lvnodal} stops at the hull, and for a reason: outside it,
interpolation becomes extrapolation, and extrapolation weights are
signed --- an extrapolated value is no longer an average of nodal
values.  With $a>0$ and a first cell whose variance decreases toward
the boundary, the linear continuation can reach zero before $y=0$
(pricing floors it there, so the continued curve kinks rather than
turning negative), and a rising continuation can exceed $v_{\rm hi}$
freely.  Every guarantee in this chapter that leans on the box is
therefore scoped to the hull and to the flat-clamped wings.  The wing
is a stated convention --- the same standing as the tail
conventions of Chapters~2--3 --- and the chapter's contract table
records it.
\end{remark}

\begin{remark}[Lee's bound, inherited]
\label{rem:lvlee}
A bounded field bounds the smile, and the bound is worth deriving
because Chapter~3 had to legislate what comes out here for free.  The
first step is that the marched call is \emph{monotone in the field}:
raising local variance anywhere can only raise call prices.  To see
it, let $c_{v}$ and $c_{v+\delta v}$ solve \eqref{eq:lvdupire} for
fields $v$ and $v+\delta v$ with $\delta v\ge0$, and subtract the two
equations.  Their difference $s=c_{v+\delta v}-c_{v}$ starts at
$s(0,\cdot)=0$ and satisfies
\[
\partial_\tau s
=\tfrac12\,(v+\delta v)\,y^{2}\,\partial_{yy}s
+\tfrac12\,\delta v\,y^{2}\,\partial_{yy}c_{v}
\]
(add and subtract $\tfrac12(v+\delta v)y^{2}\partial_{yy}c_{v}$ to
check).  The source term is nonnegative --- $\delta v\ge0$ by
assumption, $\partial_{yy}c_{v}\ge0$ by \cref{prop:lvpropagate} ---
and a heat equation driven by a nonnegative source from a zero start
stays nonnegative (the maximum principle again).  So $s\ge0$.  Now
compare with the constant field $v_{\rm hi}$.  With flat-clamped
wings, $v\le v_{\rm hi}$ holds on the entire strike line, so every
call price is dominated by its price in the Black--Scholes model with
variance rate $v_{\rm hi}$; and since the Black price increases in
total variance (vega is positive), the domination transfers to implied
variance:
\[
w(k,\tau)\;\le\;v_{\rm hi}\,\tau\qquad\text{for every }k .
\]
Total implied variance bounded in strike means its wing slope decays
to zero: the surface sits strictly inside Lee's cap of $2$
automatically, where Chapter~3 had to impose a cap by hand
(\cref{sec:svilee}).  A steep released continuation ($a$ free) escapes
$v_{\rm hi}$, and with it this bound.
\end{remark}

\section{Pricing on the lattice}
\label{sec:lvlattice}

\subsection{The implicit march and its M-matrix}

The forward equation must now be solved numerically, and this section
chooses the scheme with one criterion above all others: it must
inherit \cref{prop:lvpropagate}, not merely approximate it.

First the setting.  The equation is solved on a strike grid over
$[0,y_{\max}]$, with the two end values pinned: $c(\cdot,0)=1$ (a call
struck at zero delivers the forward, worth $1$ in normalized units)
and $c(\cdot,y_{\max})=0$ (far enough out, the call is worthless).
Pinned end values are called \emph{Dirichlet} boundaries.  The march
then steps maturity forward in increments $\Delta\tau$, and at each
step it must decide at which time level to evaluate the strike
derivative.  The \emph{fully implicit} choice --- evaluate at the new
level --- is the one this chapter defends.

On the grid, the second derivative becomes a weighted second
difference, and the weights are worth deriving because their signs
carry the whole theory.  The grid need not be uniform; write
$\Delta y_i^-=y_i-y_{i-1}$ and $\Delta y_i^+=y_{i+1}-y_i$ for the gaps
around node $i$.  Taylor-expand the two neighbors of a smooth function
$c$ around $y_i$:
\[
c_{i\pm1}=c_i\pm\Delta y_i^{\pm}\,c'(y_i)
+\tfrac12\,(\Delta y_i^{\pm})^{2}\,c''(y_i)+O(\Delta y^{3}).
\]
Multiply the ``$+$'' expansion by $\Delta y_i^-$, the ``$-$''
expansion by $\Delta y_i^+$, and add: the first-derivative terms
cancel, leaving
\[
\Delta y_i^-\,c_{i+1}+\Delta y_i^+\,c_{i-1}
-(\Delta y_i^-+\Delta y_i^+)\,c_i
=\tfrac12\,\Delta y_i^+\Delta y_i^-(\Delta y_i^++\Delta y_i^-)\,
c''(y_i)+O(\Delta y^{4}),
\]
and solving for $c''$ gives the three-point weights.  Multiplying by
$\tfrac12 v_i y_i^{2}$, the discrete generator $\mathcal L$ encoding
$\tfrac12 v y^{2}\partial_{yy}$ has the stencil
\begin{equation}
\mathcal L_{i,i\mp1}
=\frac{v_i\,y_i^{2}}{(\Delta y_i^-+\Delta y_i^+)\,\Delta y_i^{\mp}}\;\ge0,
\qquad
\mathcal L_{ii}=-\big(\mathcal L_{i,i-1}+\mathcal L_{i,i+1}\big),
\label{eq:lvstencil}
\end{equation}
which reduces to the familiar $(1,-2,1)/\Delta y^{2}$ pattern on a
uniform grid.  Note the two signs that will do all the work: the
off-diagonal entries are nonnegative, and each diagonal is exactly
minus the sum of its row's off-diagonals.

An implicit step then solves a linear system for the vector $U^{n+1}$
of calls at the new maturity, the operator evaluated at the new level:
\begin{equation}
\mathcal M^{n+1}U^{n+1}=U^{n}+(\text{boundary terms}),
\qquad
\mathcal M^{n+1}=I-\Delta\tau\,\mathcal L^{n+1}.
\label{eq:lvstep}
\end{equation}
Because the stencil couples only neighbors, $\mathcal M$ is
tridiagonal, and each step is one tridiagonal solve --- linear cost in
the number of nodes.  The boundary terms carry the pinned values $1$
and $0$ through the stencil's end columns, and two later proofs lean
on them, so write them out once as a vector $b^{n+1}$: only the row
whose lower neighbor is the boundary node $y=0$ has a nonzero entry,
\[
b^{n+1}_i=\Delta\tau\,\mathcal L^{n+1}_{i,i-1}\times1
\quad(\text{the pinned call value at }y=0),
\qquad
b^{n+1}_i=0\ \text{ at every other row}
\]
--- the far boundary would contribute through its own stencil arm, but
its pinned value is $0$.

Now the promised discipline.  \Cref{prop:lvpropagate} was proved for
the continuous equation.  Does the discrete march \emph{inherit} it
--- no invented negative densities, no broken calendar order --- or
merely approximate it, with errors of unknown sign?  The answer runs
through one classical notion from linear algebra.  A matrix with
positive diagonal, nonpositive off-diagonal entries and enough
diagonal dominance is called an \emph{M-matrix}, and the property that
sign pattern buys is exactly the one needed here: the inverse of an
M-matrix has no negative entries.

\begin{proposition}[The step matrix is an M-matrix]
\label{prop:lvmmatrix}
For any $v\ge0$ and any $\Delta\tau>0$, the matrix
$\mathcal M=I-\Delta\tau\mathcal L$ of \eqref{eq:lvstep} has positive
diagonal, non-positive off-diagonals, and each diagonal entry exceeds the
absolute sum of its row's off-diagonals by at least $1$ (exactly $1$ at
fully interior rows).  Hence $\mathcal M$ is a strictly diagonally
dominant nonsingular M-matrix, and $\mathcal M^{-1}\ge0$ elementwise.
\end{proposition}

\begin{proof}
The sign pattern reads off \eqref{eq:lvstencil}: the off-diagonals of
$\mathcal M$ are $-\Delta\tau\mathcal L_{i,i\mp1}\le0$; the diagonal is
$1+\Delta\tau(\mathcal L_{i,i-1}+\mathcal L_{i,i+1})\ge1$; and the
dominance gap is the $1$ contributed by the identity at interior rows,
more at rows whose stencil arm reaches a boundary column (that
neighbor is not among the unknowns).  This is the defining pattern of
a nonsingular M-matrix, and such matrices are inverse-positive
\citep{BermanPlemmons1994}.  A self-contained argument fits in four
lines.  To solve $\mathcal M u=b$ with $b\ge0$, split
$\mathcal M=D-\mathcal N$ into its diagonal $D>0$ and off-diagonal
part $\mathcal N\ge0$, and iterate Jacobi's method from zero:
\[
u^{(m+1)}=D^{-1}\big(\mathcal N u^{(m)}+b\big),\qquad u^{(0)}=0 .
\]
Every iterate is nonnegative, being assembled from nonnegative pieces;
the iteration converges, because every row of $D^{-1}\mathcal N$ sums
to less than one (that is diagonal dominance), so each update
multiplies the largest error entry by a factor below one; and the
limit solves $\mathcal M u=b$.  So
$b\ge0$ forces $u\ge0$, for every such $b$ --- which is the statement
$\mathcal M^{-1}\ge0$, column by column.
\end{proof}

Inverse positivity is the discrete counterpart of heat flow never
turning a nonnegative temperature profile negative, and it translates
directly into the property numerical
analysts call \emph{monotonicity} --- order in, order out --- and into
stability in the strongest everyday sense: the largest absolute value
in the solution can never grow.

\begin{corollary}[Discrete maximum principle, unconditionally]
\label{cor:lvdmp}
For any $\Delta\tau>0$ the scheme \eqref{eq:lvstep} is \emph{monotone} ---
$U^{n}\ge V^{n}$ with the same boundary data implies $U^{n+1}\ge V^{n+1}$
--- and every step keeps the values between the extremes of the previous
step and the boundary data: the march is unconditionally
$\ell^{\infty}$-stable (the maximum absolute value never grows), with
no step-size restriction, and cannot create
new extrema.
\end{corollary}

\begin{proof}
Monotonicity is inverse positivity applied to a difference: if
$U^{n}\ge V^{n}$, the two systems \eqref{eq:lvstep} subtract to
$\mathcal M^{n+1}(U^{n+1}-V^{n+1})=U^{n}-V^{n}\ge0$, and multiplying
by $\mathcal M^{-1}\ge0$ preserves the sign.  For the upper bound, let
$U_\ast$ be the larger of $\max_iU^{n}_i$ and the boundary values $1$
and $0$, and test the constant vector $U_\ast\1$, where $\1$ is the
all-ones vector (a vector here, not the indicator of
\cref{prop:lvforward}'s proof).  At a fully interior row, the row of
$\mathcal M$ sums to $1$ (because the stencil's row sums to zero), so
$(\mathcal M\,U_\ast\1)_i=U_\ast\ge U^{n}_i$.  At a row next to a
boundary, the arm that would reach the boundary column is missing from
$\mathcal M$, leaving the larger value
$(\mathcal M\,U_\ast\1)_i=U_\ast+\Delta\tau\,\mathcal L_{i,\rm bd}\,U_\ast$,
with $\mathcal L_{i,\rm bd}$ the missing arm's weight; the right-hand
side there is $U^{n}_i+b^{n+1}_i$ with
$b^{n+1}_i\le\Delta\tau\,\mathcal L_{i,\rm bd}\times1$.  Since
$U_\ast$ dominates both $U^{n}_i$ and the boundary value $1$, the left
side dominates the right at every row:
$\mathcal M(U_\ast\1)\ge U^{n}+b^{n+1}$ entrywise.  Multiplying by
$\mathcal M^{-1}\ge0$ gives $U_\ast\1\ge U^{n+1}$.  The lower bound is
symmetric.
\end{proof}

\begin{remark}[What the M-matrix argument does and does not prove]
\label{rem:lvscope}
\Cref{prop:lvmmatrix,cor:lvdmp} control the call \emph{values}: order
preservation, boundedness, no invented extrema.  They do not by
themselves give nonnegative second differences --- a bounded monotone
vector can still be locally concave --- so ``the scheme cannot
manufacture negative densities'' would need a separate
convexity-preservation proof for the actual stencil and boundary
treatment, and this chapter does not claim one.  What it does instead
is measure: the scheme is monotone and stable by theorem, and the
butterfly and calendar cleanliness of the \emph{output} is checked on
the price grid with every fit.  \Cref{sec:lvexamples} reports the
measured minima for the chapter's two calibrated sheets; they sit at
the scale of solver rounding.
\end{remark}

\subsection{Why not Crank--Nicolson}

The implicit march is only first-order accurate in the time step, and
the standard second-order upgrade is Crank--Nicolson: evaluate the
operator half at the old level and half at the new, replacing
\eqref{eq:lvstep} by
\[
\big(I-\tfrac{\Delta\tau}{2}\mathcal L\big)U^{n+1}
=\big(I+\tfrac{\Delta\tau}{2}\mathcal L\big)U^{n}.
\]
Look at what the upgrade costs.  The implicit factor on the left is
still an M-matrix --- that part survives.  But monotonicity now also
needs the \emph{explicit} factor on the right to map nonnegative
vectors to nonnegative vectors, and its diagonal entries are
$1-\tfrac{\Delta\tau}{2}(\mathcal L_{i,i-1}+\mathcal L_{i,i+1})$:
nonnegative only when the step is small enough,
\begin{equation}
\Delta\tau\;\le\;\frac{2}{\mathcal L_{i,i-1}+\mathcal L_{i,i+1}}
\;=\;\frac{2\,\Delta y^{2}}{v_i\,y_i^{2}}
\quad\text{(uniform grid)}.
\label{eq:lvcfl}
\end{equation}
This is a step-size restriction of the Courant--Friedrichs--Lewy (CFL)
type --- the generic fate of schemes with an explicit part --- and it
tightens with the \emph{square} of
the strike step.  Put numbers in: at a $40\%$ local volatility on a
$\Delta y=0.005$ lattice --- the resolution \cref{sec:lvinfo} will
show short-dated smiles require --- the bound is
$\MacLvLatCflBound$ years, more than thirty times smaller than the
standard $0.01$ marching step.  Monotone Crank--Nicolson on realistic
lattices is simply not on offer.  When the bound is violated, the
scheme remains second-order \emph{accurate} but is no longer
\emph{monotone}, and the payoff kink at $y=1$ seeds the oscillation
that \cref{fig:lvmonotone} displays.

\begin{figure}[htbp]
\centering
\paperfig{0.9\textwidth}{fig_lv_monotone.pdf}
\caption{The same flat $40\%$ surface, the same $\Delta y=0.005$
lattice, marched to $\tau=0.1$ in ten steps by the same solver ---
only the scheme differs.  (a)~Implicit Euler: the discrete density is
a clean bell, minimum \MacLvLatImplMinDens\ in the plotted window.
(b)~Undamped Crank--Nicolson: the kink oscillation swings the discrete
density to \MacLvLatCnMinDens.}
\label{fig:lvmonotone}
\end{figure}

\Cref{fig:lvmonotone} shows what breaking monotonicity does to the
object we actually care about: the discrete density, i.e.\ the second
differences of the marched prices.  Both panels march the same flat
$40\%$ surface on the same $\Delta y=0.005$ lattice to $\tau=0.1$ in
ten steps; only the scheme differs.  In panel~(a), implicit Euler: the
second differences form a clean bell --- a discrete Gaussian, as they
should --- whose minimum over the plotted window is
\MacLvLatImplMinDens: nonnegative, exactly as \cref{cor:lvdmp}
guarantees.  In panel~(b), undamped Crank--Nicolson: the oscillation
seeded by the payoff kink swings the discrete density down to
\MacLvLatCnMinDens\ --- a negative density, which is butterfly
arbitrage, manufactured by the scheme that is, in the textbook
ordering, the more accurate of the two.  Damped start-up steps --- a
few implicit steps before switching to Crank--Nicolson --- are the
standard repair \citep{GilesCarter2006}.  The reference implementation
declines the trade altogether: it pays one order of time accuracy and
keeps \cref{cor:lvdmp} unconditionally.

\subsection{Reading out, and the operator's blind spot}

The observation operator is plain: the normalized call at a quote's
$(\tau,y)$ is read off the marched solution by interpolation in
strike, and quoted expiries are forced to be lattice time-nodes, so no
interpolation in maturity is ever needed.

One principle from inverse problems completes the section, because it
dictates part of the chapter's protocol.  An optimizer facing a
\emph{fixed} discretization will happily bend its parameters to cancel
the operator's own discretization error.  The residual measured on the
fitting operator therefore flatters the fit, and validating a
reconstruction with the operator that produced it has a name in the
inverse-problems literature: the \emph{inverse crime}.  The chapter's
protocol answers it by repricing every fit on a \emph{refined}
operator --- half the strike step, four times the time steps --- and
reporting both numbers.  The gap between them is not noise; it is the
measured size of the compensation, and \cref{sec:lvexamples} will show
it is the dominant error term at the short end.

\section{What the quotes determine}
\label{sec:lvinfo}

Before writing an objective function, a chapter should ask whether the
problem it is about to pose has one answer.  This one does not, in three
distinct ways --- one about resolution, one about averages, one about
whole regions of the field --- and each dictates a piece of the
calibration's design.

\paragraph{The length scale.}
Everything interesting about a smile at maturity $\tau$ happens within a
few multiples of the diffusive width $\sigma\sqrt{\tau}$ of the money.
This single quantity dictates where both grids of
\cref{sec:lvsheet,sec:lvlattice} must put their points.  \emph{Vertex
placement}: strike columns are spaced by delta --- a fixed ladder of
exercise probabilities mapped through $\PhiN^{-1}$, rather than fixed
steps of strike --- so parameter density follows where optionality
lives.  \emph{Vertex coverage}: an axis sized by
the longest expiry undersamples a short one, whose entire smile can fall
between two columns --- a two-day smile at $20\%$ volatility is about
$1.5\%$ wide, an order of magnitude narrower than a one-year smile --- so
every expiry must be guaranteed several columns \emph{within its own}
range, spread evenly across it.  \emph{Lattice steps}: the same width
bounds the strike step from above, and if the \emph{quote} spacing is
finer than the lattice step, the fitted smile oscillates at quote
frequency until the lattice out-resolves it.  A grid that does not resolve
the data's length scale fits an alias, not the data; the exact placement
rules used in this chapter are collected in \cref{app:lvnumerics}.

\paragraph{An averaging fact, and the floor.}
At the money, implied and local variance are related as average to
integrand.  If the field near the money depends only on time,
$v(\tau,y)\approx v(\tau,1)$, the model is Black--Scholes with a
time-dependent volatility and \emph{exactly}
$w(0,\tau)=\int_0^{\tau}v(s,1)\,\dd s$; in general this holds to leading
order near the money.  An average can only be matched by an integrand that
crosses it: whenever the ATM term structure is \emph{increasing}, early
local variance must sit strictly \emph{below} the shortest implied
variance.  A parameter floor set carelessly --- say at a round $5\%$
volatility --- therefore makes some low-volatility short-dated term
structures literally unfittable, with the optimizer pinned to the box and
the misfit looking like a mystery.  The box of \cref{prop:lvnodal} must be
generous and \emph{adaptive}: a cap a fixed multiple of the largest
implied volatility in view, a floor a fraction of the \emph{smallest ATM}
implied --- keyed to ATM quotes deliberately, so that one noisy deep-wing
print cannot lower the floor where it does real stabilizing work.

\paragraph{Rows the quotes cannot see.}
The quotes at the first expiry $\tau_1$ pin --- through the averaging fact
--- only the \emph{integral} of local variance over $[0,\tau_1]$.  Any
vertex row strictly inside that interval is individually unidentified: the
optimizer can ring one row up and the next down at no cost in fit, and
under noise it will (this is precisely how the rippled fit of
\cref{fig:lvwrongway} spent its freedom at the shortest expiry, wandering
\MacLvWrongEdgePts\ volatility points at a span edge held by a single
quote).  The cure is not more data but a declared prior: tie each sub-front
row to the first identified row --- the \emph{front tie} of
\cref{sec:lvcalib} --- turning an unidentified subspace into a stated
constant continuation.

\paragraph{Two sheets, one quote set.}
The three paragraphs above are instances of one geometry, and it is better
seen than told.  A fit like this chapter's carries a couple of hundred
nodal parameters, and even when the quotes outnumber them several-fold
--- nearly five to one on the frozen SPY chain of \cref{sec:lvexamples},
between two and three on NVDA --- the quotes cluster near the money of a
handful of expiries, and a vanilla constrains local variance only through
a smoothing integral along diffusion paths.  Whole regions of the sheet
(deep wings, sub-front rows, times beyond the last expiry) are therefore
weakly determined no matter how the counting goes.  \Cref{fig:lvidentify} makes it concrete on the synthetic
running example: pushing the sheet's unquoted deep-put column (at
$y=\MacLvIdentColY$, below every quoted strike) down by
\MacLvIdentMovePts\ full volatility points changes the reprice of the
entire quote set by at most \MacLvIdentMaxDiffBp\ bp ---
\MacLvIdentRmsDiffBp\ bp rms --- concentrated in the longest expiry's
lowest strikes, the only quotes whose diffusion reaches the moved column
at all.  A thousand basis points of surface, twenty of quotes.

\begin{figure}[htbp]
\centering
\paperfig{\textwidth}{fig_lv_identify.pdf}
\caption{What sparse quotes do not see.  (a)~The calibrated synthetic
sheet at $\tau=0.5$ (solid) and the same sheet with its unquoted deep-put
column pushed down by \MacLvIdentMovePts\ volatility points (dashed); the
shaded band is the quoted region, where the two are indistinguishable.
(b)~The absolute reprice change at every quote of every expiry: at most
\MacLvIdentMaxDiffBp\ bp, and only where the longest expiry's lowest
strikes feel the moved column.  Low quote error is \emph{not} surface
recovery.}
\label{fig:lvidentify}
\end{figure}

The consequence is philosophical and operational at once.  What a
calibration returns is one \emph{representative} of an equivalence class
of sheets the data cannot distinguish, and every device in the next
section's objective is a declared choice of representative: the roughness
term selects the smoothest member near a reference, the box clips the
unidentified extremes, the front tie pins the rows below the first expiry,
and the seed decides which basin a non-convex search lands in.  None of
this is hidden tuning; it is the ``convention'' column of the book's
thesis, operating inside a calibration.  The worked example of
\cref{sec:lvexamples} accordingly reports quote error and surface error as
two separate numbers.

\section{The objective and its derivatives}
\label{sec:lvcalib}

\subsection{The objective}

The pieces now assemble into a bound-constrained least-squares problem
over the nodal variances:
\begin{equation}
\min_{\theta\in[v_{\rm lo},\,v_{\rm hi}]^{n_\theta}}
\;\tfrac12\Big(
\sum_{q}r_q(\theta)^{2}
\;+\;\lambda\,\big\|\Gamma(\theta-\theta_{\rm ref})\big\|^{2}
\;+\;\lambda_0\,\big\|\theta_{1,\cdot}-\theta_{2,\cdot}\big\|^{2}
\Big).
\label{eq:lvobjective}
\end{equation}
Its three blocks, in order: a data term, a roughness penalty, and a
tie between the first two vertex rows --- $\theta_{1,\cdot}$ is the
$\tau=0$ row, $\theta_{2,\cdot}$ the row at the first quoted expiry.
Each block was argued for before it was written down; we
take them in turn.

The \emph{data block} compares the marched call $c_\theta$ at quote
$q$'s coordinates with the quoted one, in volatility units:
$r_q=\big(c_\theta(\tau_q,y_q)-c_q\big)/\eta_q$, where $\eta_q$ is one
volatility percent of the quote's vega, floored.  Recall what vega is:
the derivative of the Black price with respect to implied volatility
--- the exchange rate between price units and volatility units.
Dividing each price error by it makes every residual read as a
volatility error, so a deep wing's tiny price errors and the money's
large ones compete on the same scale; the floor (the same device as
Chapter~2's vega floor $\eta$; this chapter's constant is in
\cref{app:lvnumerics}) keeps the scale from vanishing where vega does.
This is the residual design of \cref{sec:calibration} --- solve in
price space, measure in volatility space --- carried over unchanged,
and when band edges rather than mids are trusted, the band objective
of that section --- residuals measured to the nearest bid--ask band
edge, zero inside the band --- applies verbatim.

The \emph{roughness block} $\lambda\|\Gamma(\theta-\theta_{\rm ref})\|^{2}$
is Tikhonov regularization --- the standard name for adding a penalty
so that, among sheets fitting the data equally well, the optimizer
returns the smoothest.  The operator $\Gamma$ takes second divided
differences along strike within each maturity row and along time
within each strike column --- the discrete curvature of the sheet ---
normalized by the local vertex spacing, so that a tightly spaced money
column is not penalized more than a wide wing one (on a uniform grid
it reduces to the $(1,-2,1)$ stencil).  The penalty pulls toward a
reference sheet $\theta_{\rm ref}$ --- flat, at the median ATM
variance, in this chapter's protocol --- and it is the declared answer
to \cref{sec:lvinfo}: of the equivalence class the data allows, return
the smoothest member near the reference.

The \emph{front tie}, at weight $\lambda_0$, pins the $\tau=0$ row to
the first quoted row, closing the invisible sub-front subspace of
\cref{sec:lvinfo}.

Two further row types plug into the same stack when their instruments
are present, and are only named here.  A variance-swap quote adds one
scalar residual per quoted expiry, constraining the deep-put tail
integral that vanillas cannot see --- Chapter~5 develops the integral,
and it is exactly the instrument that justifies releasing the wing
slope $a$ of \cref{rem:lvhull}.  And prior structures from a previous
fit enter as reference terms in the same
$\Gamma$-and-$\theta_{\rm ref}$ position --- Chapter~10 develops when
and how much to trust them.

\subsection{Every derivative from the same factorization}

A least-squares solver runs on the Jacobian: the matrix of derivatives
of every residual with respect to every parameter.  For a model priced
by a PDE march, the naive route is finite differences --- perturb
$\theta_\ell$, re-march the surface, subtract --- which costs
$n_\theta$ extra marches per iteration, hundreds here.  The sheet's
structure gives something better: \emph{exact} derivatives, for the
price of extra right-hand sides through a factorization already
computed.  The key is linearity.  By \eqref{eq:lvp1} the field depends
linearly on $\theta$, so the discrete generator does too, and its
derivative in one parameter is itself a stencil:
\[
\frac{\partial\mathcal L}{\partial\theta_\ell}
=\mathcal L[\mathcal H_\ell],
\]
the stencil \eqref{eq:lvstencil} with the field $v$ replaced by the
single hat function $\mathcal H_\ell$ --- zero outside vertex $\ell$'s
star of triangles.

\begin{proposition}[Tangent: same factor, extra right-hand sides]
\label{prop:lvtangent}
The sensitivities $s^{n}_\ell=\partial U^{n}/\partial\theta_\ell$ obey
\begin{equation}
\mathcal M^{n+1}s^{n+1}_\ell
=s^{n}_\ell+\Delta\tau\,\mathcal L[\mathcal H_\ell]\,U^{n+1},
\qquad s^{0}_\ell=0,
\label{eq:lvtangent}
\end{equation}
where $\mathcal L[\mathcal H_\ell]\,U^{n+1}$ applies the boundary-inclusive
stencil to the solution extended by its fixed boundary values.  This is the
discretization of the continuous sensitivity equation
$\partial_\tau s_\ell=\tfrac12 v y^{2}\partial_{yy}s_\ell
+\tfrac12\mathcal H_\ell\,y^{2}\partial_{yy}c$.
\end{proposition}

\begin{proof}
Write one step over the interior unknowns as
$\mathcal M^{n+1}U^{n+1}=U^{n}+b^{n+1}$, with $b^{n+1}$ the boundary
vector displayed after \eqref{eq:lvstep} --- it carries the
Dirichlet values through the stencil's boundary columns, and
therefore depends on $\theta$ through the field at boundary-adjacent
nodes.  Differentiate the whole equation with respect to
$\theta_\ell$, using the product rule on the left:
\[
\Big(\partial_{\theta_\ell}\mathcal M^{n+1}\Big)U^{n+1}
+\mathcal M^{n+1}s^{n+1}_\ell
=s^{n}_\ell+\partial_{\theta_\ell}b^{n+1}.
\]
By linearity,
$\partial_{\theta_\ell}\mathcal M^{n+1}
=-\Delta\tau\,\mathcal L[\mathcal H_\ell]$ exactly --- no
approximation is involved.  Move it to the right-hand side:
\[
\mathcal M^{n+1}s^{n+1}_\ell
=s^{n}_\ell
+\Delta\tau\,\mathcal L[\mathcal H_\ell]\,U^{n+1}
+\partial_{\theta_\ell}b^{n+1},
\]
and the last two terms combine into the boundary-inclusive product of
the statement, because the boundary \emph{values} themselves ($1$ and
$0$) do not move with $\theta$ --- only the stencil weights through
which they enter do.  The start is $s^{0}_\ell=0$: the payoff does not
depend on the field.
\end{proof}

Read \eqref{eq:lvtangent} twice.  As a computation: the sensitivity
system has the \emph{same} matrix $\mathcal M^{n+1}$ as the value
system --- only the right-hand side differs --- so every sensitivity
marches behind the same tridiagonal factorization as the value solve.
The entire $n_\theta$-column Jacobian costs one factorization plus
$n_\theta$ extra back-substitutions per step; the columns form a block
of right-hand sides that a solver processes together.  As probability:
the source term says that vertex $\ell$ influences a price exactly
through the diffusion's visits to $\ell$'s star of triangles, weighted
by the convexity $\partial_{yy}c$ it finds there.

\begin{figure}[htbp]
\centering
\paperfig{\textwidth}{fig_lv_influence.pdf}
\caption{The tangent system, seen and audited.  (a)~Sensitivity of the
priced call curve to the single vertex at $(\tau,y)=(0.25,1)$, across
the four quoted expiries.  (b)~Every analytic column of the Jacobian
against a central finite difference of the full march, over all
\MacLvInflVtx\ vertices: worst relative disagreement
\MacLvInflMaxRel.}
\label{fig:lvinfluence}
\end{figure}

\Cref{fig:lvinfluence} shows both readings on the synthetic market.
In panel~(a), the sensitivity of the priced call curve to the single
vertex at $(\tau,y)=(0.25,1)$, plotted across the four quoted
expiries.  At $\tau=0.10$ the curve is identically zero: an earlier
expiry cannot feel a later vertex, because the march has not reached
it yet.  At $\tau=0.25$, the vertex's own maturity, the sensitivity is
a sharp cone centered on its strike.  At the two later expiries the
cone spreads and flattens: the diffusion gradually forgets where it
collected its variance.  In panel~(b), the audit: every analytic
Jacobian column is compared against a central finite difference of the
full march, over all \MacLvInflVtx\ vertices, and the worst relative
disagreement is \MacLvInflMaxRel.  The derivative is exact, not
approximate.

The optimization itself is then standard machinery used plainly: a
trust-region reflective least-squares iteration on
\eqref{eq:lvobjective} --- a method of the Gauss--Newton family
(linearize the residuals at the current iterate, solve the resulting
linear least squares, repeat) that handles
box bounds natively, shrinking its steps rather than leaving the box
\citep{BranchColemanLi1999} --- with the analytic Jacobian of
\cref{prop:lvtangent} supplied at every iterate and a flat seed at the
reference sheet.  The tangent route's Jacobian cost grows linearly
with $n_\theta$; the classical \emph{adjoint} alternative --- run the
march backward once through its transposed steps, collecting the whole
gradient in a single sweep, at a cost independent of the
parameter count --- is derived in \cref{app:lvnumerics} for the reader
who will need it at larger grids.

\section{Examples and the contract}
\label{sec:lvexamples}

All fits in this section use one protocol --- the grids of
\cref{sec:lvinfo}, the objective of \cref{sec:lvcalib}, every constant
recorded in \cref{app:lvnumerics} --- applied unchanged to a synthetic
market whose answer is known and to the book's frozen SPY and NVDA
chains.  No fit is tuned separately.

\subsection{The round trip, and its two separate numbers}

The synthetic market prices a known smooth field through an
\emph{independent} dense scheme (so the exercise cannot flatter the
chapter's march by validating it against itself), quotes it at four
expiries, and hands the \MacLvRtNQuotes\ clean quotes to the calibration
from a flat seed.  Per \cref{sec:lvinfo}, the result is two numbers, not
one.  The \emph{quotes} reprice to \MacLvRtRmsErrBp\,bp rms and
\MacLvRtMaxErrBp\,bp at worst --- the fit is essentially exact in the
space the optimizer sees.  Separately, the \emph{field itself} agrees with
the truth to \MacLvRtSurfRmsPts\ volatility points rms on the
quote-covered region, with a maximum of \MacLvRtSurfMaxPts\ points ---
more than a hundred times the quote error, concentrated exactly where
\cref{fig:lvidentify} said the quotes do not look: between expiries and
near the span edges (\cref{fig:lvrecovery}b).  Both numbers are honest;
only reporting the first would not be.

\begin{figure}[htbp]
\centering
\paperfig{\textwidth}{fig_lv_recovery.pdf}
\caption{The synthetic round trip.  (a)~Target quotes (dots) and the
recovered fit (lines) at the four expiries: rms \MacLvRtRmsErrBp\,bp, max
\MacLvRtMaxErrBp\,bp.  (b)~The second number: $|$fitted $-$ true$|$ local
volatility over the quote-covered region, quote positions dotted.  The
error lives between the expiries and at the edges of the quoted spans ---
largest where the nearest quote is farthest.}
\label{fig:lvrecovery}
\end{figure}

\subsection{The frozen chains}

The real test fits one sheet per underlier to every prepared mid quote of
the snapshot's eight expiries at once: \MacLvSpyQuotes\ quotes against
\MacLvSpyVtx\ nodal variances for SPY, \MacLvNvdaQuotes\ against
\MacLvNvdaVtx\ for NVDA.  \Cref{fig:lvsheet} already showed the SPY sheet
itself; \cref{fig:lvfit} shows four of its expiries against the quotes,
priced by the \emph{one} calibrated diffusion --- a two-day smile and a
year-and-a-half LEAP from the same list of positive numbers.

\begin{figure}[htbp]
\centering
\paperfig{\textwidth}{fig_lv_fit.pdf}
\caption{Four SPY expiries priced by the single calibrated sheet, against
the frozen quotes (dots: mids; whiskers: bid--ask).  Panel titles give the
expiry and its distance; annotations the per-expiry rms on the fitting
operator.}
\label{fig:lvfit}
\end{figure}

\begin{figure}[htbp]
\centering
\paperfig{\textwidth}{fig_lv_rms.pdf}
\caption{Per-expiry rms for SPY (a) and NVDA (b), on the fitting operator
(grey) and on the refined operator (blue), log scale.  The refined bars
are the ones to read: the gap between the pairs is the compensated
discretization error of \cref{sec:lvlattice}, and it owns the short end.}
\label{fig:lvrms}
\end{figure}

\Cref{fig:lvrms} reports the chapter's two metrics side by side.  On the
fitting operator the surfaces read \MacLvSpyRmsBp\,bp (SPY) and
\MacLvNvdaRmsBp\,bp (NVDA) rms across all quotes --- NVDA's wider spreads
and harder short-dated smiles carry more residual everywhere.  Repriced on
the refined operator the same surfaces read \MacLvSpyConvBp\ and
\MacLvNvdaConvBp\,bp, and the worst single expiry --- in both names the
one-day front --- reads \MacLvSpyWorstConvBp\ and
\MacLvNvdaWorstConvBp\,bp.  The gap between the paired bars is the point,
and it is the operator's blind spot of \cref{sec:lvlattice} made
quantitative: the optimizer partially cancels the coarse march's error at
the expiries it is fitted on, the in-operator residual flatters
accordingly, and the number to quote is the refined one.  Where the two
agree (the long end) the fit is genuinely converged; where they diverge
(day-scale maturities) the remaining error is discretization, not
disagreement with the quotes.  The measured cleanliness of
\cref{rem:lvscope} completes the report: across every expiry of both
calibrated sheets, the minimum divided second difference of the marched
prices is \MacLvSpyButterflyMin\ (SPY) and \MacLvNvdaButterflyMin\ (NVDA),
and the minimum adjacent-expiry increment is \MacLvSpyCalendarMin\ and
\MacLvNvdaCalendarMin\ --- solver rounding, on both counts, as
\cref{prop:lvpropagate} predicts and \cref{rem:lvscope} declines to assume.

\subsection{The contract}
\label{sec:lvcontract}

\begin{table}[htbp]
\centering\small
\begin{tabular}{@{}>{\raggedright\arraybackslash}p{0.31\textwidth}
>{\raggedright\arraybackslash}p{0.31\textwidth}
>{\raggedright\arraybackslash}p{0.30\textwidth}@{}}
\toprule
\textbf{Structural} & \textbf{Numerically controlled} &
\textbf{Outside the claim}\\
\midrule
Box bounds imply a positive field on the hull; a positive field implies
butterfly-free, calendar-ordered marginals (\cref{prop:lvnodal,%
prop:lvpropagate}); every optimizer iterate is an admissible diffusion &
Butterfly, calendar and bound cleanliness of the marched prices, measured
per fit; refined-operator reprice beside the in-operator residual &
Convexity preservation of the discrete scheme as a theorem; the
extrapolation region beyond the hull, where basis weights are signed\\[4pt]
Monotone, unconditionally $\ell^\infty$-stable march
(\cref{prop:lvmmatrix,cor:lvdmp}); exact tangent Jacobian from the same
factorization (\cref{prop:lvtangent}) & Tangent columns audited against
finite differences; an independent discretization cross-checks the march &
Uniqueness of the calibrated sheet: the data determines an equivalence
class, the conventions of \cref{sec:lvcalib} a representative\\[4pt]
Bounded field $\Rightarrow$ implied variance bounded in strike, strictly
inside Lee's cap on the hull (\cref{rem:lvlee}) & Grid placement rules
resolving each expiry's $\sigma\sqrt\tau$; adaptive variance box keyed to
ATM quotes & Jumps and scheduled-event spikes (the clock, Chapter~8);
spot--volatility dynamics (Chapter~9); the deep-put tail's shape without a
variance-swap quote (Chapter~5)\\
\bottomrule
\end{tabular}
\caption{The local-volatility model contract: theorem, numerical control,
and excluded scope.}
\label{tab:lvcontract}
\end{table}

The chapter's separation thesis reads off the table.  \emph{Validity} is
structural and cheap: a box.  \emph{Information} is scarcer than the
formula \eqref{eq:lvextract} pretends: quotes determine the field only
through smoothing integrals, and only where they exist.  \emph{Convention}
fills the rest --- roughness, reference, tie, wing, seed.  The calibration
does not eliminate those choices; it declares them, measures what the
discretization leaks, and reports the error twice: once as the optimizer
saw it, once as it is.


\begin{chapterappendices}
\section{Numerical facts and the chapter protocol}
\label{app:lvnumerics}

\subsection{The protocol}
\label{app:lvprotocol}

Every fit in this chapter --- synthetic and real --- uses the following
recipe; no constant is tuned per name.

\emph{Vertex grid.}  Maturity rows at $\tau=0$ and at every quoted expiry.
Strike columns delta-spaced: the ladder of put deltas
$\{1,2,5,10,25,40\}\%$, the same ladder mirrored on the call side, and
the money, mapped to
strikes through
$y=\exp\big(\sigma_A\sqrt{\tau_{\max}}\;\PhiN^{-1}(\cdot)\big)$ with
$\sigma_A$ the longest expiry's ATM volatility; the two end columns are
extended to cover the pooled quoted range, so no quote falls outside the
hull.  Coverage rule: every expiry must own at least $6$ columns inside
its own traded range (\cref{sec:lvinfo}), enforced as follows: list
the columns falling inside the expiry's range together with the
range's two endpoints, and while the widest gap between consecutive
entries exceeds one fifth of the range, split that gap with a new
column; this is what densifies the short-dated money in
\cref{fig:lvsheet}.

\emph{Lattice.}  Strike step
\[
\Delta y=\min\big(0.01,\;0.15\,\min_e\sigma_e\sqrt{\tau_e}\big),
\]
capped at $800$ nodes over $[0,\,\max(2.5,\,1.05\,y_{n_y})]$ with
$y_{n_y}$ the highest vertex column, and $y=1$ a node; every quoted
expiry is a time node and each maturity interval receives
$\max(8,\lceil\Delta\tau_{\rm int}/0.01\rceil)$ implicit steps.
The refined (audit) operator halves the strike step and quadruples each
interval's step count.

\emph{Box and objective.}  Adaptive box
\[
v_{\rm lo}=\min\big(0.0025,\,(\tfrac12\min_e\sigma_{{\rm ATM},e})^{2}\big),
\qquad
v_{\rm hi}=\min\big(\max(0.36,\,(3\,\sigma_{\max})^{2}),\,16\big)
\]
--- the floor a fraction of the smallest ATM implied per the averaging
fact, the cap a multiple of the largest implied in view.  Seed and reference
$\theta^{0}=\theta_{\rm ref}$ flat at the median ATM variance.  Weights
$\lambda=10^{-2}$, $\lambda_0=1$; residual scale $\eta_q$ equal to one
volatility percent of the quote's Black vega, floored at $10^{-3}$.  The
snapshot carries no variance-swap quotes, so the wing multiple stays at
its default $a=0$ (flat clamp) throughout.

\emph{Solver.}  Trust-region reflective least squares
\citep{BranchColemanLi1999} with the analytic Jacobian of
\cref{prop:lvtangent}, termination tolerances $10^{-8}$, at most $200$
objective evaluations; the fits of \cref{sec:lvexamples} converged in
\MacLvSpyEvals\ (SPY), \MacLvNvdaEvals\ (NVDA) and \MacLvRtEvals\
(synthetic) evaluations from their flat seeds.

\subsection{The synthetic market}
\label{app:lvsynthetic}

The running example's truth is the smooth skewed field
\[
\sigma_{\rm loc}(\tau,y)
=0.21-0.10\,\tanh\!\Big(\frac{y-1}{0.22}\Big)
+0.03\,e^{-2\tau},
\qquad v=\sigma_{\rm loc}^{2},
\]
priced by an \emph{independently implemented} dense march (uniform
$\Delta y=0.0025$ on $[0,3]$, $\Delta\tau=0.002$, fully implicit) --- a
second discretization on different grids, so the round trip of
\cref{fig:lvrecovery} cross-checks the chapter's march rather than
validating it against itself.  Quotes are read at the four expiries
$\tau\in\{0.10,0.25,0.50,1.00\}$, at $31$ strikes per expiry uniform in
the standardized moneyness $\zeta\in[-2.2,2.2]$ of \cref{sec:mcs}, with
$\sigma_{\rm ref}=0.24$, i.e.\ $k=0.24\sqrt{\tau}\,\zeta$ (the coarse arm
of \cref{fig:lvwrongway}a uses $13$).  The ``noise'' of
\cref{sec:lvillposed} is a deterministic alternating
$\pm50$\,bp volatility ripple --- no randomness anywhere in the chapter's
figures.  The synthetic sheet has rows $\{0\}\cup\{\text{expiries}\}$ and
the delta-spaced columns at scale $0.24$ plus one deliberate deep-put
column at $y=\MacLvIdentColY$ for the experiment of
\cref{fig:lvidentify}; its box is $[0.005,0.36]$.

\subsection{The adjoint}
\label{app:lvadjoint}

The tangent route of \cref{prop:lvtangent} prices the Jacobian at
$O(N_\tau N_y\,n_\theta)$ per evaluation, for a march of $N_\tau$ steps on
$N_y$ lattice nodes.  The classical alternative makes the cost of a
\emph{gradient} independent of $n_\theta$; deriving it costs a few lines,
so the chapter records it for the grids on which it will one day pay.
The derivation follows the standard optimal-control pattern ---
adjoin the marching constraints with multiplier vectors $p^{n}$, the
\emph{adjoint states}, require stationarity, and read the result as a
backward march; a reader meeting adjoints for the first time can treat
\eqref{eq:lvadjoint}--\eqref{eq:lvadjgrad} as a recipe and check it
numerically against the tangent Jacobian.

\begin{proposition}[Adjoint: one backward sweep]
\label{prop:lvadjoint}
Let $\mathcal J(\theta)$ be a scalar objective depending on the march
through the observed values $\mathcal O_nU^{n}$, with $\mathcal O_n$ fixed
observation rows.  Define adjoint states by the backward recursion
\begin{equation}
(\mathcal M^{n})^{\!\top}p^{n}
=\mathcal O_n^{\top}\,
\frac{\partial\mathcal J}{\partial(\mathcal O_nU^{n})}+p^{\,n+1},
\qquad p^{\,N+1}=0 .
\label{eq:lvadjoint}
\end{equation}
Then
\begin{equation}
\frac{\partial\mathcal J}{\partial\theta_\ell}
=\sum_{n}\Delta\tau\,(p^{n})^{\!\top}\,
\mathcal L[\mathcal H_\ell]\,U^{n},
\label{eq:lvadjgrad}
\end{equation}
with the boundary-inclusive product of \cref{prop:lvtangent}, at the cost
of one backward solve, $O(N_\tau N_y)$, plus the $O(n_\theta)$
accumulation of \eqref{eq:lvadjgrad}.
\end{proposition}

\begin{proof}
Form the Lagrangian (fraktur $\mathfrak L$, kept visually apart from
the stencil operator $\mathcal L$)
$\mathfrak{L}=\mathcal J-\sum_n (p^{n})^{\top}
\big(\mathcal M^{n}U^{n}-U^{n-1}-b^{n}\big)$, with $b^{n}$ the boundary
terms of \cref{prop:lvtangent}.  Stationarity in each $U^{n}$ --- which
appears in step $n$'s system and in step $n{+}1$'s right-hand side ---
gives \eqref{eq:lvadjoint}; the $\theta$-derivative of $\mathfrak{L}$ then
retains the explicit dependence of $\mathcal M^{n}$ and $b^{n}$, which
combine as in \cref{prop:lvtangent} into the boundary-inclusive
\eqref{eq:lvadjgrad}.  Structurally this is summation by parts: the
transpose of a lower block-bidiagonal system is upper block-bidiagonal,
i.e.\ a backward march with transposed tridiagonal factors.
\end{proof}

\subsection{Measured cleanliness}
\label{app:lvdiag}

The diagnostics quoted in \cref{sec:lvexamples} are computed on the
fitting operator's price grid: the butterfly proxy is the minimum divided
second difference $\big(U_{i+1}-U_i\big)/\Delta y_i^{+}
-\big(U_i-U_{i-1}\big)/\Delta y_i^{-}$ over all interior nodes and all
expiries; the calendar proxy is the minimum of
$U^{(e+1)}-U^{(e)}$ over adjacent expiries at fixed $y$ (the lattice
prices are already in the normalized units of \cref{sec:calendar}, so the
fixed-$y$ comparison is the right one); and the bounds check verifies
$0\le U\le1$ to the same tolerance.  For both calibrated sheets all three
sit at solver rounding (\cref{sec:lvexamples}); they are measured per fit
precisely because \cref{rem:lvscope} declines to promote them to a
theorem.

\end{chapterappendices}
